% =====================================================================
% P1 Appendix A.4 --- Theorem 3 (branch impossibility).
%
% Source of truth: drafts/proofs-props8-12.md v1.4, Sections 1--4
% (FS-8, FS-9, FS-9.1).  This file is a fragment: it is meant to be
% \input under \appendix \section{Proofs} after A.3.
%
% Theorem-like environments use the explicit paper numbering
% (Theorems 3A--3C, Lemmas 1--2, Remark 1.1).  Equation tags
% (1.1)--(1.4), (2.1)--(2.4), (3.1)--(3.2), (4.1)--(4.3) follow the
% frozen source; do not renumber without updating the main text.
% =====================================================================

\newtheorem*{aIVthmA}{Theorem 3A}
\newtheorem*{aIVthmB}{Theorem 3B}
\newtheorem*{aIVthmC}{Theorem 3C}
\newtheorem*{aIVlemA}{Lemma 1}
\newtheorem*{aIVlemB}{Lemma 2}
\newtheorem*{aIVremA}{Remark 1.1}

% Fixed equation tags: keep the frozen printed numbers while giving
% hyperref a unique destination for every display.
\providecommand{\theHequation}{}
\let\aIVtheHequation\theHequation
\newcounter{aIVanchor}
\renewcommand{\theHequation}{aIV.\arabic{aIVanchor}}
\newcommand{\aIVfixedtag}[1]{\stepcounter{aIVanchor}\tag{#1}}

\subsection{Branch impossibility (Theorem~3)}\label{app:a4}

This appendix gives the finite-sample statements and proofs behind
Theorem~3.  The population versions are the fixed-threshold statements
summarized in Section~5.1; the statements below are stronger in that
they allow an arbitrary finite sample and a data-dependent threshold.  Theorem~3A is the finite-sample branch lower bound,
Theorem~3B the non-identification result, and Theorem~3C the sharp
coupling bound together with its minimax consequence.  Theorem~3 in the
main text is the union of 3A--3C.

\paragraph{Setting.}
Let $D_n=\{(X_i,Y_i)\}_{i=1}^n$ be a calibration sample and let
$(X_{n+1},Y_{n+1},B_{n+1})$ be a test point, where $B\in\{0,1\}$ is a
post-decision branch (for example, rebalance versus hold).  On the
coverage side, the branch labels $B_i$ may be observed (branch-wise
calibration) or unavailable (branch-blind rules).  A scalar conformity
score is $s(x,y)$; for CQR one may take
$s(x,y)=\max(\hat q_{\mathrm{lo}}(x)-y,\,y-\hat q_{\mathrm{hi}}(x))$.
A \emph{branch-blind score rule} is
\[
C_n(x)=\{y:s(x,y)\le\hat q\},
\]
where $\hat q=\hat q(D_n)$ depends only on the calibration pairs
$(X_i,Y_i)$ and on independent randomization, and not on any $B_i$.
For a fixed rule (conditionally on $D_n$ and the randomization), let
$F(x)=\{y:y\notin C_n(x)\}$ be the failure set.

Define the conditional quantities
\[
\begin{aligned}
\alpha(x)&=\Pr(Y\in F(x)\mid X=x), &
\alpha&=\Pr(Y\in F(X))=\mathbb E[\alpha(X)],\\
q_b(x)&=\Pr(B=b\mid X=x), &
q_b&=\Pr(B=b)=\mathbb E[q_b(X)],\\
f_b(x)&=\Pr(Y\in F(x),\,B=b\mid X=x).
\end{aligned}
\]
Then the branch-$b$ coverage is
\begin{equation}\label{eq:a4-1.1}
\operatorname{cov}_b(C)
=\Pr(Y\in C(X)\mid B=b)
=1-\frac{\mathbb E[f_b(X)]}{q_b},
\qquad q_b>0.
\aIVfixedtag{1.1}
\end{equation}
When $q_b=0$ the conditioning event has probability zero and branch $b$
is omitted; all statements below are for $q_b>0$.

\subsubsection{Auxiliary lemmas}

\begin{aIVlemA}[feasible interval]
For every $x$,
\begin{equation}\label{eq:a4-1.2}
\max\bigl(0,\alpha(x)+q_b(x)-1\bigr)
\le f_b(x)\le
\min\bigl(\alpha(x),q_b(x)\bigr).
\aIVfixedtag{1.2}
\end{equation}
Conversely, every value in this interval is realized by some joint law
of $(X,Y,B)$ with the given law of $(X,Y)$ and the given $q_b(x)$.
When $0<\alpha(x)<1$ and $0<q_b(x)<1$ the interval has nonempty
interior; its length is
$\min(\alpha,q_b)-\max(0,\alpha+q_b-1)>0$.  If
$\alpha(x)\in\{0,1\}$ or $q_b(x)\in\{0,1\}$ the interval may degenerate
to a point; such cases are handled by continuity and do not change the
coverage lower bounds.
\end{aIVlemA}

\begin{proof}
For the upper bound, $f_b(x)\le\Pr(F(x)\mid X=x)=\alpha(x)$ and
$f_b(x)\le\Pr(B=b\mid X=x)=q_b(x)$.  For the lower bound, $f_b(x)\ge0$
and
\[
\Pr(F(x)^c,B=b\mid X=x)=q_b(x)-f_b(x)
\le\Pr(F(x)^c\mid X=x)=1-\alpha(x),
\]
hence $f_b(x)\ge\alpha(x)+q_b(x)-1$.

For the converse, put mass $f_b(x)$ of $B=b$ on $F(x)$ and mass
$q_b(x)-f_b(x)$ on $F(x)^c$.  The two nonnegativity conditions are
exactly the two sides of \eqref{eq:a4-1.2}.  The length claim follows by
considering $\alpha+q_b\le1$ and $\alpha+q_b>1$ separately.
\end{proof}

\begin{aIVlemB}[tail-function coupling]
Fix $x$.  Let $S=s(X,Y)$ and $\alpha_x(q)=\Pr(S>q\mid X=x)$.  Define
the nonincreasing functions
\begin{equation}\label{eq:a4-1.3}
F_b(q)=\min\bigl(\alpha_x(q),q_b(x)\bigr),
\qquad
F'_b(q)=\max\bigl(0,\alpha_x(q)+q_b(x)-1\bigr).
\aIVfixedtag{1.3}
\end{equation}
Then there exist conditional laws $P_x$ and $P'_x$ of $B\mid(X=x,S)$
such that
\begin{equation}\label{eq:a4-1.4}
\begin{aligned}
P_x(B=b,\,S>q\mid X=x)&=F_b(q),\\
P'_x(B=b,\,S>q\mid X=x)&=F'_b(q),
\end{aligned}
\aIVfixedtag{1.4}
\end{equation}
and $P_x(B=b\mid X=x)=P'_x(B=b\mid X=x)=q_b(x)$.  Moreover
$F'_b(q)\le F_b(q)$, and the inequality is strict whenever
$0<\alpha_x(q)<1$ and $0<q_b(x)<1$.
\end{aIVlemB}

\begin{proof}
Both $F_b$ and $F'_b$ are right-continuous and nonincreasing in $q$,
with limits
\[
F_b(-\infty)=F'_b(-\infty)=q_b(x),
\qquad
F_b(\infty)=F'_b(\infty)=0,
\]
and $0\le F'_b(q)\le F_b(q)\le\alpha_x(q)$.

By the Stieltjes measure theorem, there is a unique finite Borel measure
$\mu_b$ on $\mathbb R$.  It satisfies
$\mu_b((q,\infty))=F_b(q)$ for every $q$.  Similarly,
\[
G_b(q):=\alpha_x(q)-F_b(q)=\max\bigl(\alpha_x(q)-q_b(x),0\bigr)
\]
is right-continuous and nonincreasing, with
$G_b(-\infty)=1-q_b(x)$ and $G_b(\infty)=0$, and
$0\le G_b\le\alpha_x$; let $\mu_{1-b}$ be the associated finite measure.
For every half-open interval $(a,b]$ with
$-\infty\le a<b\le\infty$,
\[
\mu_b((a,b])+\mu_{1-b}((a,b])
=F_b(a)-F_b(b)+G_b(a)-G_b(b)
=\alpha_x(a)-\alpha_x(b)
=\mu_S((a,b]).
\]
The finite disjoint unions of such intervals form an algebra that
generates the Borel $\sigma$-field, and the total masses agree:
\[
\mu_b(\mathbb R)+\mu_{1-b}(\mathbb R)
=q_b(x)+(1-q_b(x))=1=\mu_S(\mathbb R).
\]
Dynkin's uniqueness theorem therefore gives
$\mu_b+\mu_{1-b}=\mu_S$ on all Borel sets.  In particular
$0\le\mu_b\le\mu_S$, so the Radon--Nikodym derivative
$d\mu_b/d\mu_S$ takes values in $[0,1]$.  Setting
\[
\Pr(B=b\mid X=x,S=s)=\frac{d\mu_b}{d\mu_S}(s)
\]
defines the kernel $P_x$, and
$P_x(B=b,S>q\mid X=x)=\mu_b((q,\infty))=F_b(q)$.
The same construction with $F'_b$, whose complement
$\alpha_x-F'_b=\min(\alpha_x,1-q_b(x))$ is also nonincreasing, gives
$P'_x$.  Measurability in $x$ follows by defining the tail functions on
rational $q$ and extending by right-continuity.

Finally, $F'_b\le F_b$ because
$\alpha+q_b-1\le\min(\alpha,q_b)$ and $0\le\min(\alpha,q_b)$.  If
$0<\alpha_x(q)<1$ and $0<q_b(x)<1$, then $F'_b<F_b$: when
$\alpha_x+q_b\le1$ we have $F'_b=0<F_b$, and when
$\alpha_x+q_b>1$ we have
$F'_b=\alpha_x+q_b-1<\min(\alpha_x,q_b)=F_b$ because $\alpha_x<1$ and
$q_b<1$.
\end{proof}

\begin{aIVremA}[data-dependent thresholds]
Lemma~2 defines the conditional law of $B\mid(X,S)$ simultaneously for
all thresholds $q$.  Consequently, if $\hat q$ is a function of the
calibration sample and the rule is branch-blind, the law of $\hat q$ is
the same under the two couplings of Lemma~2; one simply conditions on
$D_n$ and evaluates the construction at $q=\hat q$.  This is the point
that lets Theorem~3B cover data-dependent thresholds.
\end{aIVremA}

\subsubsection{Finite-sample branch lower bound (Theorem~3A)}

\begin{aIVthmA}
Let $C$ be any prediction set, possibly depending on the calibration
sample (including its branch labels), and suppose its marginal
miscoverage satisfies $\Pr(Y\notin C(X))\le\alpha$.  Then for every
branch $b$ with $q_b>0$,
\begin{equation}\label{eq:a4-2.1}
\operatorname{cov}_b(C)
=\Pr(Y\in C(X)\mid B=b)
\ge1-\frac{\alpha}{q_b}.
\aIVfixedtag{2.1}
\end{equation}
If in addition $C$ is branch-blind, then under calibration--test
independence (in particular under i.i.d.\ sampling), conditionally on
$D_n$ and the randomization,
\begin{equation}\label{eq:a4-2.2}
\operatorname{cov}_b(C\mid D_n)
\ge1-\frac{\mathbb E[\min(\alpha(X),q_b(X))\mid D_n]}{q_b}
\ge1-\frac{\alpha(D_n)}{q_b},
\aIVfixedtag{2.2}
\end{equation}
where $\alpha(D_n)=\mathbb E[\alpha(X)\mid D_n]$.  Taking expectations
over $D_n$ and using $\alpha=\mathbb E[\alpha(D_n)]$ recovers
\eqref{eq:a4-2.1}.  If $q_b=0$, the conditioning event is undefined and
the branch is omitted.
\end{aIVthmA}

\begin{proof}
The law of total probability gives
\begin{equation}\label{eq:a4-2.3}
\begin{aligned}
\Pr(Y\notin C(X))
&=q_b\Pr(Y\notin C(X)\mid B=b)
 +(1-q_b)\Pr(Y\notin C(X)\mid B=1-b)\\
&\ge q_b\Pr(Y\notin C(X)\mid B=b).
\end{aligned}
\aIVfixedtag{2.3}
\end{equation}
The left side is at most $\alpha$, so
$\Pr(Y\notin C(X)\mid B=b)\le\alpha/q_b$, which is
\eqref{eq:a4-2.1}.

For branch-blind $C$, condition on $D_n$ and the independent
randomization.  Then $C$ is fixed, and Lemma~1 applied to $f_b(X)$ gives
$f_b(X)\le\min(\alpha(X),q_b(X))$.  Hence
\begin{equation}\label{eq:a4-2.4}
\operatorname{cov}_b(C\mid D_n)
=1-\frac{\mathbb E[f_b(X)\mid D_n]}{q_b}
\ge1-\frac{\mathbb E[\min(\alpha(X),q_b(X))\mid D_n]}{q_b}.
\aIVfixedtag{2.4}
\end{equation}
Since $\min(\alpha,q_b)\le\alpha$, the right side is at least
$1-\alpha(D_n)/q_b$.  Taking expectations over $D_n$ and using
$\mathbb E[\alpha(D_n)]=\alpha$ gives
$\operatorname{cov}_b(C)\ge1-\alpha/q_b$, which is again
\eqref{eq:a4-2.1}.
\end{proof}

\paragraph{Equality and tightness.}
\begin{itemize}
\item Equality in \eqref{eq:a4-2.1} holds if and only if all marginal
failure mass falls in branch $b$, that is, $f_b(X)=\alpha(X)$ almost
surely (for $q_b>0$).
\item Equality in \eqref{eq:a4-2.4} holds if and only if
$f_b(X)=\min(\alpha(X),q_b(X))$ almost surely, that is, branch $b$ takes
the full feasible upper mass.
\item If $\alpha(X)\le q_b(X)$ almost surely, the two bounds in
\eqref{eq:a4-2.2} coincide, and a coupling that places all failures in
branch $b$ attains $\operatorname{cov}_b=1-\alpha/q_b$.
\item If $q_b(X)<\alpha(X)$ on a set of positive probability, then
$\mathbb E[\min(\alpha(X),q_b(X))]<\mathbb E[\alpha(X)]=\alpha$, so the
sharp bound in \eqref{eq:a4-2.2} is strictly larger than
$1-\alpha/q_b$.  Branch $b$ is too small to absorb all failures, and the
marginal bound is loose.  This observation is the starting point of
Theorem~3C.
\end{itemize}

\paragraph{Where branch-blindness enters.}
The pointwise inequality
\[
f_b(x)\le\min(\alpha(x),q_b(x))
\]
from Lemma~1 holds for every rule, including rules that use calibration or
test-time branch labels.  Branch-blindness is used only to make the
conditional allocation $B\mid(X,S)$ invariant to the rule, so that the
right side of \eqref{eq:a4-2.2} can serve as a sharp bound common to all
compatible couplings and can be used in the attainability construction of
Theorem~3C.  Without branch-blindness, two cases arise.  A rule that uses
only calibration branch labels still satisfies \eqref{eq:a4-2.2}, but its
right side is no longer a common sharp bound, because the threshold law
changes with the coupling.  A rule that uses the test-time branch label
also satisfies the pointwise inequality and \eqref{eq:a4-2.2}, but its
allocation may change with the coupling, so the same common sharpness and
attainability conclusions are not asserted.

\subsubsection{Non-identification from pre-decision data (Theorem~3B)}

\begin{aIVthmB}
Fix the marginal law of $X$, the joint law of $(X,Y)$ (hence the score
law), and branch probabilities $q_b(x)\in(0,1)$.  Let $C_n$ be any
branch-blind score rule, with threshold $\hat q=\hat q(D_n)$ possibly
data-dependent.  Then there exist two joint laws $P$ and $P'$ on
$(X,Y,B)$ such that:
\begin{enumerate}
\item[(i)] $(X,Y)$ has the same law under $P$ and $P'$; consequently the
law of the calibration sample's $(X,Y)$ part, the law of $\hat q$, and
the output law of $C_n$ are the same under $P$ and $P'$;
\item[(ii)] $P(B=b\mid X=x)=P'(B=b\mid X=x)=q_b(x)$ for every $x$;
\item[(iii)] for every $n\ge1$, if
\[
\Pr_{D_n}\Bigl(\Pr_X\bigl(0<\alpha_{D_n}(X)<1\bigr)>0\Bigr)>0,
\]
then the branch-$b$ coverages are strictly different:
\begin{equation}\label{eq:a4-3.1}
\begin{aligned}
&\bigl|\operatorname{cov}_b^P(C_n)
-\operatorname{cov}_b^{P'}(C_n)\bigr|\\
&\qquad
=\frac{\mathbb E\bigl[\min(\alpha(X),q_b(X))
-\max(0,\alpha(X)+q_b(X)-1)\bigr]}{q_b}
>0.
\end{aligned}
\aIVfixedtag{3.1}
\end{equation}
\end{enumerate}
Here $\alpha_{D_n}(x)=\Pr(Y\notin C_n(x)\mid X=x,D_n)$; $\Pr_X$ is over
$X$, $\Pr_{D_n}$ is over $D_n$, and the expectation in
\eqref{eq:a4-3.1} is over $(X,D_n)$.  If
$\alpha_{D_n}(X)\in\{0,1\}$ almost surely, the two couplings may give
the same coverage; see the degenerate case below.
\end{aIVthmB}

\paragraph{Construction.}
For each $x$, apply Lemma~2 with $S=s(X,Y)$ and
$\alpha_x(q)=\Pr(S>q\mid X=x)$ to define $B\mid(X=x,S)$ by
\begin{equation}\label{eq:a4-3.2}
\begin{aligned}
P:\quad & \Pr(B=b,\,S>q\mid X=x)
=\min\bigl(\alpha_x(q),q_b(x)\bigr),\\
P':\quad & \Pr(B=b,\,S>q\mid X=x)
=\max\bigl(0,\alpha_x(q)+q_b(x)-1\bigr).
\end{aligned}
\aIVfixedtag{3.2}
\end{equation}
Combining these conditional laws with the fixed marginal law of
$(X,Y)$ gives $P$ and $P'$.  The construction is simultaneous in $q$
and therefore does not depend on $\hat q$.

\paragraph{Verification.}
\begin{enumerate}
\item \emph{The $(X,Y)$ law is unchanged.}  The construction changes
only $B\mid(X,Y)$, not $(X,Y)$.  Hence the calibration sample's
$(X,Y)$ part has the same law, $\hat q$ has the same law, and, since
$C_n$ does not read $B$, the output of $C_n$ is pointwise the same under
$P$ and $P'$.

\item \emph{The branch probabilities are unchanged.}  Let
$q\to-\infty$ in \eqref{eq:a4-3.2}.  Then $\alpha_x(q)\to1$, so both
$\min(1,q_b(x))$ and $\max(0,1+q_b(x)-1)$ equal $q_b(x)$.  Hence
$P(B=b\mid X=x)=P'(B=b\mid X=x)=q_b(x)$.

\item \emph{The coverage gap.}  Condition on $D_n$, so that $\hat q$ is
fixed, and write $\alpha(x)=\alpha_x(\hat q)$.  Under $P$, the failure
mass of branch $b$ is $\min(\alpha(x),q_b(x))$; under $P'$, it is
$\max(0,\alpha(x)+q_b(x)-1)$.  By \eqref{eq:a4-1.1}, the absolute
coverage gap is \eqref{eq:a4-3.1}.  Lemma~2 gives nonnegativity; the
integrand is strictly positive on
$\{0<\alpha(x)<1,\;0<q_b(x)<1\}$, so the expectation is strictly
positive whenever that set has positive measure.  The $P$ side has the
larger failure mass and therefore the lower coverage; the $P'$ side has
the smaller failure mass and the higher coverage.

\item \emph{Every $n$ is covered.}  The construction is a single i.i.d.\
joint law, so it holds for every sample size; the rule may use any finite
$D_n$, whose distribution is the same under the two laws.
\end{enumerate}

\paragraph{Degenerate rules.}
If $\alpha(X)\in\{0,1\}$ almost surely, the failure set is determined by
$X$, and $f_b$ is uniquely determined by $q_b$ and the failure
indicator.  The two couplings then give the same coverage.  In that case
branch coverage is identifiable from the $(X,Y)$ law, but the rule is
degenerate: at each $x$ it either covers or fails with probability one,
which is not the typical split-conformal rule at level $1-\alpha$.
Nondegenerate continuous scores, whose conditional support covers the
threshold $\hat q$, normally give $\alpha(X)\in(0,1)$ almost surely.  The
main text states the nondegenerate version; to include degenerate rules,
add independent randomization that flips coverage with arbitrarily small
probability, so that $\alpha(X)\in(0,1)$ and the conclusion is unchanged.

\paragraph{Relation to the population statement.}
The population proposition fixes a threshold $q$ and a cell-wise
$\alpha$, and constructs $P,P'$ piecewise.  Theorem~3B writes the same
idea as a tail-function coupling that holds simultaneously for all $q$;
hence no simplification ``$\hat q$ is a fixed $q$'' is needed, no
expectation over $\hat q$ has to be taken separately, and the conclusion
holds for every $n$ with the explicit gap \eqref{eq:a4-3.1}.  The
piecewise formula is recovered by setting
$D=\min(\alpha,q_b)-\max(0,\alpha+q_b-1)$ and splitting into four cells:
\[
\begin{array}{ll}
\alpha+q_b\le1,\ q_b\le\alpha: & D/q_b=1,\\
\alpha+q_b\le1,\ q_b>\alpha: & D/q_b=\alpha/q_b,\\
\alpha+q_b>1,\ q_b\ge\alpha: & D/q_b=(1-q_b)/q_b,\\
\alpha+q_b>1,\ q_b<\alpha: & D/q_b=(1-\alpha)/q_b.
\end{array}
\]
The boundary between the last two cells is $q_b=\alpha$, not
$q_b=1-\alpha$.  When $\alpha>1/2$, every
$q_b\in(1-\alpha,\alpha)$ satisfies $\alpha+q_b>1$ and $q_b<\alpha$, so
it belongs to the fourth cell; classifying only by $q_b>1-\alpha$ would
put these points in the third cell.

\subsubsection{Sharp coupling bound and minimax (Theorem~3C)}

\begin{aIVthmC}
Fix the law of $(X,Y)$, the branch probabilities $q_b(x)$, and a
branch-blind score rule $C_n$.  Let $\mathcal P$ be the class of all
joint laws of $(X,Y,B)$ with the given $(X,Y)$ law and the given
$q_b(x)$.  Then, for $q_b>0$,
\begin{equation}\label{eq:a4-4.1}
\inf_{P\in\mathcal P}\operatorname{cov}_b^P(C_n)
=1-\frac{\mathbb E[\min(\alpha(X),q_b(X))]}{q_b},
\aIVfixedtag{4.1}
\end{equation}
where $\alpha(x)=\Pr(Y\notin C_n(x)\mid X=x,D_n)$ is first defined
conditionally on $D_n$ and then averaged over $D_n$ and $X$.  This
quantity is the same under every $P\in\mathcal P$, because the $(X,Y)$
law is fixed and the rule is branch-blind.  The infimum is attained by
the $P$ side of \eqref{eq:a4-3.2}; the $P'$ side gives the other
endpoint, namely the supremum
$1-\mathbb E[\max(0,\alpha(X)+q_b(X)-1)]/q_b$.
\end{aIVthmC}

\begin{proof}
Theorem~3A, inequality \eqref{eq:a4-2.4}, after averaging over $D_n$,
yields the lower bound in \eqref{eq:a4-4.1} for every
$P\in\mathcal P$; hence the infimum is at least the right side of
\eqref{eq:a4-4.1}.  Conversely, the $P$ side of Lemma~2 defines
$B\mid(X,S)$ for all $q$ simultaneously and attains
$f_b=\min(\alpha,q_b)$ at $q=\hat q$ for every realization of $D_n$.
It is therefore a fixed i.i.d.\ law in $\mathcal P$ whose coverage equals
the right side of \eqref{eq:a4-4.1}; hence the infimum is at most that
value.  This proves \eqref{eq:a4-4.1}.  The statement about the $P'$
side follows from the same lemma with $F'_b$ in place of $F_b$.
\end{proof}

\paragraph{The constant case.}
In the population constant case, $\alpha(x)\equiv\alpha$ and
$q_b(x)\equiv q_b$ (for example, a fixed threshold rule),
\eqref{eq:a4-4.1} becomes
\begin{equation}\label{eq:a4-4.2}
\inf_{P\in\mathcal P}\operatorname{cov}_b^P
=1-\frac{\min(\alpha,q_b)}{q_b}
=\max\Bigl(1-\frac{\alpha}{q_b},0\Bigr).
\aIVfixedtag{4.2}
\end{equation}
If $q_b\ge\alpha$, the value is $1-\alpha/q_b$ and is attained by the
$P$ side; if $q_b<\alpha$, the value is $0$ and is attained by the $P$
side as well (branch $b$ fails on all of its mass).  Combined with
\eqref{eq:a4-2.1}, for a fixed rule the infimum over compatible
couplings is exactly $\max(1-\alpha/q_b,0)$.

If one further takes the supremum over rules in the constant class
(choose the rule, then let the coupling be worst), every rule in the
class has the same worst-case value, so the minimax value in
Theorem~3(c) is exactly the trivial bound in this class, provided the
class is nonempty.  The next paragraph shows the same for the larger
class $\alpha(X)\le q_b(X)$ almost surely.

\paragraph{Cell-wise covariates.}
If $\alpha(x)$ varies with $x$, \eqref{eq:a4-4.1} is the sharp fixed-rule
bound, while the population bound $1-\alpha/q_b$ remains valid but is
generally not tight:
\begin{equation}\label{eq:a4-4.3}
1-\frac{\mathbb E[\min(\alpha(X),q_b(X))]}{q_b}
\ge1-\frac{\min(\alpha,q_b)}{q_b}
=\max\Bigl(1-\frac{\alpha}{q_b},0\Bigr).
\aIVfixedtag{4.3}
\end{equation}
The inequality uses
$\mathbb E[\min(\alpha(X),q_b(X))]\le\min(\mathbb E\alpha,\mathbb E q_b)$,
which follows from concavity of the minimum.  Equality holds, for
example, if $\alpha(X)\le q_b(X)$ almost surely or if
$q_b(X)\le\alpha(X)$ almost surely, since then the minimum is linear on
the support; other cancellations can also produce equality, so these
conditions are sufficient but not necessary.  In particular, if
$\alpha(X)\le q_b(X)$ almost surely, then
$\min(\alpha(X),q_b(X))=\alpha(X)$ and the sharp value is still
$1-\alpha/q_b$ even when $\alpha(x)$ varies.  This is why the class
$\alpha(X)\le q_b(X)$ almost surely in Theorem~3(c) is more general than
the constant class.  When both orderings occur on sets of positive
measure, the inequality in \eqref{eq:a4-4.3} is strict.

Assume the non-atomic conditional-law condition: $Y\mid X=x$ is non-atomic
for almost every $x$.  The larger class is nonempty whenever $q_b\ge\alpha$
and $q_b>0$.  Put $a(x)=(\alpha/q_b)q_b(x)$, so $0\le a(x)\le q_b(x)$ and
$\mathbb E[a(X)]=\alpha$.  For $a(x)<1$, let $F_x(y)=\Pr(Y\le y\mid X=x)$
and define the branch-blind score
\[
s(x,y)=\frac{F_x(y)}{1-a(x)},
\]
with the fixed threshold $\hat q=1$; for $a(x)=1$, set
$s(x,y)=+\infty$.  The failure event is
$\{s(x,Y)>1\}=\{F_x(Y)>1-a(x)\}$.  Because $Y\mid X=x$ is non-atomic,
$F_x(Y)\sim\mathrm{Unif}[0,1]$, so this event has probability $a(x)$.
The construction is an oracle existence witness: it uses the population
law $F_x$ and $q_b$, and is not proposed as an estimator from $D_n$.
The rule uses no branch label, so it belongs to the score-rule class of
Theorem~3(c).  Conversely, if $q_b<\alpha$, every
$\alpha(X)\le q_b(X)$ a.s.\ satisfies
$\mathbb E[\alpha(X)]\le q_b<\alpha$, so the class is empty.  Hence, under
the non-atomic conditional-law condition, the nonemptiness qualifier in
Theorem~3(c) is equivalent to $q_b\ge\alpha$.  For atomic $Y\mid X$ this
particular construction need not attain $a(x)$; no atomic-case claim is
made.

The bound $1-\alpha/q_b$ should therefore be read as a valid marginal
lower bound, not as a claim that the minimax value is always the trivial
one.  If $q_b(x)$ is small on a region that carries non-negligible
failure mass, branch $b$ cannot absorb all failures and the sharp bound
exceeds the trivial one.  When only the marginal bound is reported, it
should be described as potentially loose, especially for small empirical
branches.

\subsubsection{Scope of Theorem~3}

\begin{itemize}
\item Theorem~3 is an information limit, not an algorithmic defect.
Marginal coverage is a mixture over branches: a branch with small $q_b$
can have a large conditional failure rate while the marginal rate stays
$\alpha$.  Pre-decision Mondrian calibration cannot isolate a
post-decision branch whose label is unavailable at prediction time.
\item Branch-blindness is used for sharpness and reachability, not for
the validity of the basic bound.  The pointwise inequality
$f_b\le\min(\alpha,q_b)$ holds for every rule.  If the rule uses
calibration branch labels, \eqref{eq:a4-2.2} remains true but is not a
common sharp bound over compatible couplings.  If it uses the test-time
label, \eqref{eq:a4-2.2} also remains true, but the allocation may vary
across couplings, so the same common sharpness and reachability
conclusions do not follow.
\item The nondegenerate version is the paper's statement.  Degenerate
rules can be handled by independent randomization, as described above.
\item The theorem does not claim that the branch-union construction is
optimal, nor that the trivial bound is sharp for every covariate
distribution.  The sharp fixed-rule value is \eqref{eq:a4-4.1}; the
minimax equality is asserted only inside the stated rule class.
\end{itemize}

\let\theHequation\aIVtheHequation
